\documentclass[12pt]{article}
\usepackage{lineno}
%\usepackage{ametsoc}
\usepackage{natbib}
\usepackage{a4wide}
\usepackage{longtable}
\usepackage{amsfonts}
\usepackage{amssymb}

\usepackage{color}
\usepackage[pdftex]{graphicx}
\usepackage[pdftex,colorlinks, citecolor=black,linkcolor=black,urlcolor=black]{hyperref}
\definecolor{darkgreen}{rgb}{0.1,0.4,0.1}
\definecolor{darkblue}{rgb}{0.1,0.1,0.3}



\newif\ifdetail
\global\detailfalse
%\global\detailtrue

\def\bem#1{\colorbox{green}{\em #1}}

\newcommand{\beq}  { \begin{eqnarray} }
\newcommand{\eeq}  { \end{eqnarray}}
\newcommand{\beeq}  { \begin{eqnarray*} }
\newcommand{\eeeq}  { \end{eqnarray*}}
\newcommand{\eq }  [1]{Eq.~(\ref{#1})}  % Eq. (#1)
\newcommand{\eqs }  [1]{Eqs.~(\ref{#1})}
\newcommand{\sect}  [1]{Sec.~(\ref{#1})}  % Eq. (#1)
\newcommand{\fig}  [1]{Fig.~\ref{#1}}   % Fig. #1
\renewcommand{\v}[1]{{\mbox{\boldmath$ #1 $}}}  % Vektor durch Fettdruck
\newcommand{\m}    [1]{{\bf  #1 }}              % Matrix in sans serif
\newcommand{\gz}[1] {{{d #1} \over {dz}}}
\newcommand{\pt}[1]{{\partial_t #1}}
\newcommand{\ptt}[1]{{\partial_{tt} #1}}
\newcommand{\pttt}[1]{{\partial_{ttt} #1}}
\newcommand{\pttx}[1]{{{\partial^3 #1} \over {\partial t^2 \partial x}}}
\newcommand{\ptty}[1]{{{\partial^3 #1} \over {\partial t^2 \partial y}}}
\newcommand{\ptx}[1]{{{\partial^2 #1} \over {\partial t \partial x}}}
\newcommand{\pty}[1]{{{\partial^2 #1} \over {\partial t \partial y}}}
\newcommand{\pz}[1]{{\partial_z #1}}
\newcommand{\pzz}[1]{{\partial_{zz} #1}}
\newcommand{\px}[1]{{\partial_x #1}}
\newcommand{\py}[1]{{\partial_y #1}}
\newcommand{\pxi}[1]{{\partial_i #1}}
\newcommand{\pxj}[1]{{\partial_j #1}}
\newcommand{\pxk}[1]{{\partial_k #1}}
\newcommand{\pb}[1]{{{\partial #1} \over {\partial b}}}
%\newcommand{\pxy}[1]{{{\partial^2 #1} \over {\partial x \partial y}}}
\newcommand{\pxy}[1]{\partial_{xy} #1}
\newcommand{\pyy}[1]{\partial_{yy} #1}
\newcommand{\pxx}[1]{\partial_{xx} #1}
\newcommand{\kx}{{\bf k} \times}
\newcommand{\dt}[1]{{{d #1} \over {d t}}}
\newcommand{\Dt}[1]{{{D #1} \over {D t}}}
\newcommand{\Div}[1]{ \nabla \cdot #1 }
\renewcommand{\div}[1]{ \nabla_{yz} \cdot #1 }
\newcommand{\Bar}[1]{ \overline{ #1} }
\newcommand{\rvec} [1]{
    \raisebox{-1.5ex}{$\stackrel{\textstyle #1}{\neg}$} }% rotierter Vektor
\newcommand{\nablq}   {\rvec\nabla}                    % rotated nabla operator
\newcommand{\Vec}[1]{\left(\begin{array}{c} #1 \end{array}\right) }
%\renewcommand{\mho}{\mathcal{W}}
\newcommand{\pn}[1]{{{\partial #1} \over {\partial n}}}
\newcommand{\pnn}[1]{{{\partial #1} \over {\partial m}}}
\newcommand{\ps}[1]{{{\partial #1} \over {\partial s}}}
\newcommand{\p}{{\partial}}
\newcommand{\vn}{{\v \nabla}}
\newcommand{\lk}{\left(}
\newcommand{\rk}{\right)}

\newcommand{\tl}{\tilde}

\newcommand{\I}{\mathcal{I}}
\newcommand{\E}{\mathcal{E}}
\newcommand{\A}{\mathcal{A}}

\newcommand{\ep}{\epsilon}
\begin{document}

 \section*{The realistic non-hydrostatic model nh\_realistic}

nh\_realistic is pe\_realistic with a prognostic vertical velocity and a non-hydrostatic pressure.
This text describes only the differences to pe\_realistic.
Everything that is not mentioned here is as described in the documentation of pe\_realistic and pe\_simple.

The differences are the following.
The vertical velocity $w$ has its own momentum equation.
A three-dimensional pressure $p_{nh}$ makes the velocity free of divergence in every cell.
The horizontal coordinates must be Cartesian.

The package contains only what is new.
The grid, the masks, the tracers, the equation of state, the vertical mixing, the TKE closure, the isoneutral mixing, the surface pressure and the split-explicit external mode are imported from pe\_realistic.
The file \verb|momentum.py| holds the equation for $w$.
The file \verb|vmix.py| holds the vertical friction of $w$.
The files \verb|pressure.py| and \verb|solver.py| hold the equation for $p_{nh}$ and its solver.
The file \verb|model.py| holds the time step.

The formulation follows pyOM2 with \verb|enable_hydrostatic| off and the non-hydrostatic model of Oceananigans.
The solver for $p_{nh}$ follows Oceananigans.
The tests and their results are listed in the section on the nhr family of \verb|pyOM3/test/README_test_suite.md|.

\subsection*{Grid}

The grid is that of pe\_realistic with Cartesian coordinates.
The widths \verb|dxt|, \verb|dyt| and \verb|dzt| are in metres and may vary in all three directions.
The model stops with an error if \verb|coord_degree| is set.
On the sphere the equation for $w$ would need the metric term $-(u^2+v^2)/a$ and the Coriolis terms with $\cos \phi$.
nh\_realistic has neither of them.

The metric factors \verb|cost| and \verb|cosu| of pe\_realistic are therefore 1, and \verb|tantr| is 0.
The kernels imported from pe\_realistic still contain them.
The new kernels of nh\_realistic do not.
The horizontal area of a $T$ cell is $A^t = \Delta x^t \Delta y^t$, and its volume is $V^t = A^t \Delta z^t$.

The vertical velocity $w_k$ sits at the top face of the $T$ cell $k$.
This face is the centre of the $W$ cell $k$.
The $W$ cell $k$ reaches from the centre $z^t_k$ to the centre $z^t_{k+1}$.
Its height is $\Delta z^w_k$, and its volume is $A^t \Delta z^w_k$.

The mask \verb|maskW| of pe\_realistic is open at the surface face, since the free surface needs a vertical velocity there.
nh\_realistic adds the mask \verb|maskWc|, which is \verb|maskW| with the surface face closed.
The equation for $w$, its friction and the equation for $p_{nh}$ use \verb|maskWc|.
The surface face thus belongs to the free surface and not to the equation for $w$.

\subsection*{Pressure}

All pressures are divided by $\rho_0$, as in pe\_realistic.
The pressure is the sum
\beq
 p = p_h + p_s + p_{nh}
\eeq
of the hydrostatic pressure $p_h$, the surface pressure $p_s$ and the non-hydrostatic pressure $p_{nh}$.
$p_h$ and $p_s$ are computed as in pe\_realistic.

The momentum equations are
\beq \label{eq:mom}
 \pt{\v u} + \dots = - \vn_h \lk p_h + p_s + p_{nh} \rk ~,~~
 \pt w + \v u \cdot \vn w = - \pz{p_{nh}} + F_w
\eeq
where $\v u = (u,v)$ is the horizontal velocity, the dots stand for all other terms of pe\_realistic, and $F_w$ is the friction of $w$.

The equation for $w$ carries neither the buoyancy nor $p_h$ nor $p_s$.
The full equation would be $\pt w + \v u \cdot \vn w = - \pz p - g \rho/\rho_0 + F_w$.
$p_s$ does not depend on $z$.
$p_h$ is integrated downward from the surface with
\beq \label{eq:ph}
 p_h(k) = p_h(k+1) + \frac{g}{2 \rho_0} \lk \rho(k+1) + \rho(k) \rk \Delta z^w_k
\eeq
which is the integration of pe\_realistic.
\eq{eq:ph} gives at the $W$ point $k$ the vertical gradient
\beq
 \frac{p_h(k+1) - p_h(k)}{\Delta z^w_k} = - \frac{g}{2 \rho_0} \lk \rho(k+1) + \rho(k) \rk
\eeq
The right hand side is the buoyancy force averaged to the $W$ point.
The vertical gradient of $p_h$ and the buoyancy force therefore cancel exactly, also on a stretched grid.
Only $p_{nh}$ is left in the equation for $w$.

With $p_{nh} = 0$ and $w$ taken from continuity the model is pe\_realistic.
The attribute \verb|enable_nonhydrostatic| (default on) selects the non-hydrostatic model.
With it off the time step is that of pe\_realistic to the last bit.
This is tested.

\subsection*{Advection and friction of $w$}

The advection of $w$ is the energy-conserving second-order form of the momentum advection of pe\_realistic.
The transports through the faces of a $T$ cell are
\beq
 U_i = \Delta y^t \Delta z^t u_i ~,~~
 V_j = \Delta x^t \Delta z^t v_j ~,~~
 W_k = A^t w_k
\eeq
with the masks of the faces.
The transport through a face of the $W$ cell $k$ is the mean of the transports through the same face of the $T$ cells $k$ and $k+1$.
The fluxes of $w$ are
\beq
 f^x_i = - \frac{1}{4} \lk U_{i,k} + U_{i,k+1} \rk \lk w_i + w_{i+1} \rk ~,~~
 f^z_k = - \frac{1}{4} \lk W_k + W_{k+1} \rk \lk w_k + w_{k+1} \rk
\eeq
and $f^y$ correspondingly with $V$.
$f^z_k$ sits at the $T$ point $k+1$.
The advective tendency of $w_k$ is the divergence of these fluxes divided by the volume $A^t \Delta z^w_k$ of the $W$ cell.

The divergence of the transports through the $W$ cell $k$ is the mean of the divergences of the $T$ cells $k$ and $k+1$.
It vanishes when the velocity is free of divergence.
The advection then conserves $w^2$.
On a uniform grid it is the advection of $w$ of nh\_simple.

The equation for $w$ has no Coriolis term, as in pyOM2.

The horizontal friction of $w$ uses the coefficients \verb|Ahbi| and \verb|Ah| of the horizontal friction of $u$ and $v$.
It uses the operator at the $T$ points of pe\_realistic, since $w$ sits at the $T$ points in the horizontal.
For $w$ this operator is exact.
A flux exists only between two wet $W$ cells.

The vertical friction of $w$ has the flux at the $T$ point $k$ between the $W$ cells $k-1$ and $k$,
\beq
 F_k = A_k \frac{w_k - w_{k-1}}{\Delta z^t_k}
\eeq
The tendency of $w_k$ is $(F_{k+1} - F_k)/\Delta z^w_k$.
The roles of $\Delta z^t$ and $\Delta z^w$ are thus swapped with respect to the vertical friction of $u$ and $v$.
A flux exists only between two wet $W$ cells.

The coefficient $A_k$ follows the vertical friction of $u$ and $v$.
It is the constant \verb|A_v| with \verb|vmix_scheme| explicit or implicit.
With the TKE closure it is the viscosity of the closure averaged from the two $W$ points to the $T$ point.
The step is a forward step for \verb|vmix_scheme| explicit and an implicit step otherwise.
The implicit step is the tridiagonal solve of pe\_realistic on the $W$ grid.
The energy that the friction of $w$ removes is not given to the TKE closure.

\subsection*{Projection}

The velocities are first advanced with all terms except the gradient of $p_{nh}$.
This gives $u^*$, $v^*$ and $w^*$.
The surface pressure $p_s$ is already part of $u^*$ and $v^*$.
The vertical friction is also part of them.

The vertical velocity at the surface face is taken from continuity as in pe\_realistic,
\beq \label{eq:wsurf}
 w_{surf} = - \frac{1}{A^t} \sum_k \lk U_i - U_{i-1} + V_j - V_{j-1} \rk
\eeq
with $u^*$ and $v^*$ in the transports and the sum over the wet cells of the column.
With the free surface $w_{surf}$ is the change of the surface elevation in time.
With the rigid lid it is zero to the tolerance of the surface pressure solver.

The divergence of the $T$ cell is
\beq
 D = \frac{1}{V^t} \lk U_i - U_{i-1} + V_j - V_{j-1} + W_k - W_{k-1} \rk
\eeq
The velocity is corrected by
\beq \label{eq:proj}
 u_i = u^*_i - \Delta t \frac{p_{i+1} - p_i}{\Delta x^u_i} ~,~~
 v_j = v^*_j - \Delta t \frac{p_{j+1} - p_j}{\Delta y^u_j} ~,~~
 w_k = w^*_k - \Delta t \frac{p_{k+1} - p_k}{\Delta z^w_k}
\eeq
with $p = p_{nh}$, on the wet faces and the interior faces in $z$.
The surface face is closed in \eq{eq:proj}.
$w_{surf}$ therefore stays as it is.

The pressure $p_{nh}$ makes $D$ of the corrected velocity vanish in every wet cell.
Inserting \eq{eq:proj} into $D = 0$ gives
\beq \label{eq:pnh}
 a^x_i \lk p_{i+1} - p_i \rk - a^x_{i-1} \lk p_i - p_{i-1} \rk
 + a^y_j \lk p_{j+1} - p_j \rk - a^y_{j-1} \lk p_j - p_{j-1} \rk && \nonumber \\
 + a^z_k \lk p_{k+1} - p_k \rk - a^z_{k-1} \lk p_k - p_{k-1} \rk
 &=& \frac{V^t D^*}{\Delta t}
\eeq
with $D^*$ the divergence of $u^*$, $v^*$, $w^*$ and $w_{surf}$.
The coefficients are
\beq
 a^x_i = \frac{\Delta y^t \Delta z^t}{\Delta x^u_i} ~,~~
 a^y_j = \frac{\Delta x^t \Delta z^t}{\Delta y^u_j} ~,~~
 a^z_k = \frac{A^t}{\Delta z^w_k}
\eeq
on the wet faces and zero on the closed faces.
\eq{eq:pnh} is the volume times the divergence of the gradient of $p_{nh}$.

The coefficient of a face is the same in the equations of the two cells next to it.
The matrix of \eq{eq:pnh} is therefore symmetric, also on a stretched grid with topography.
This is tested.
Without the volume factor the matrix would not be symmetric on a stretched grid.

All faces at the boundary of the wet region are closed in \eq{eq:pnh}.
A solution exists only if the right hand side summed over the wet cells vanishes.
Summed over a column, the transports through the inner faces cancel.
What is left is the sum of the horizontal transports of the column and $A^t w_{surf}$.
With \eq{eq:wsurf} this sum is zero.
The right hand side thus sums to zero in every column, and \eq{eq:pnh} has a solution.

\eq{eq:pnh} fixes $p_{nh}$ only up to a constant.
After every iteration the mean of $p_{nh}$ over the wet cells is set to zero, as in Oceananigans.
The same is done with the residual.

\subsection*{Solver}

\eq{eq:pnh} is solved with the method of conjugate gradients with a preconditioner.
The solver is written for JAX and MPI in the same way as the solver for $p_s$ of pe\_simple.
The first guess is $p_{nh}$ of the previous time step.

The preconditioner solves the vertical part of \eq{eq:pnh} exactly in every column.
It drops the horizontal coefficients $a^x$ and $a^y$.
The vertical part is the tridiagonal system
\beq
 - a^z_{k-1} z_{k-1} + (1 + \epsilon) \lk a^z_k + a^z_{k-1} \rk z_k - a^z_k z_{k+1} = r_k
\eeq
for the preconditioned residual $z$ with the residual $r$.
It is solved with the Thomas algorithm.
The factor $1+\epsilon$ with $\epsilon = 0.01$ makes the system of each column regular, as in Oceananigans.
In a dry cell and in a wet cell without a wet cell above or below, the preconditioner is the identity.
The preconditioner is that of Marshall et al.\ (1997).

The preconditioner suits grids with $\Delta z \ll \Delta x$.
There the vertical coefficients $a^z$ are much larger than the horizontal ones.
In the test configuration they are larger by a factor of $3 \cdot 10^4$ to $2 \cdot 10^7$.
The preconditioner reduces the number of iterations there from 1000 to 3500 to 50 to 150 per time step.
For $\Delta z \approx \Delta x$ it helps less.

The residual $r$ is the right hand side minus the left hand side of \eq{eq:pnh}.
The divergence of the velocity corrected with the current $p_{nh}$ is $\Delta t\, r/V^t$.
The iteration stops when $\Delta t \max |r|/V^t$ is smaller than \verb|epsilon_nh| (default $10^{-12}\,$s$^{-1}$).
\verb|epsilon_nh| is thus the largest divergence that the velocity keeps.
The iteration also stops after \verb|max_iterations_nh| iterations (default \verb|max_iterations|).

The stopping condition is met by the vertical structure of $p_{nh}$.
The horizontal structure of $p_{nh}$ is less well determined when $\Delta z \ll \Delta x$.
It still changes with the round-off of the sums in the solver.
Runs with different domain decompositions therefore differ by about $10^{-8}$ relative with the default \verb|epsilon_nh|.
With \verb|epsilon_nh| $= 10^{-16}$ they differ by about $10^{-10}$.

\subsection*{Tracers}

The tracers are advected with $u$, $v$ and $w$ of the end of the previous time step, as in pe\_realistic.
In pe\_realistic $w$ follows from continuity and the velocity is free of divergence to round-off.
In nh\_realistic the velocity is free of divergence to \verb|epsilon_nh|.
A uniform tracer therefore stays uniform only to \verb|epsilon_nh| $\Delta t$ per time step.
The content of a tracer is still conserved to round-off, since the advection is in flux form.
The WENO scheme of pe\_realistic works in nh\_realistic unchanged.
For a uniform tracer it gives the same tendency as the mean of the two cells, the tracer times the divergence that the solve leaves.

\subsection*{Time step}

The order of the calls in a time step of the split-explicit scheme is
\begin{enumerate}
\item diffusivities of the TKE closure (\verb|set_tke_diffusivities|)
\item momentum tendencies, Coriolis force and hydrostatic pressure (\verb|momentum_tendency_expl|)
\item advective tendency of $w$ (\verb|w_tendency|)
\item forcing of the configuration (\verb|set_forcing|)
\item horizontal friction of $w$ and Adams--Bashforth update of $w$
\item bottom friction
\item external mode with sub cycling (\verb|solve_expl_surface|)
\item vertical friction of $u$, $v$ and of $w$
\item surface vertical velocity and the solve for $p_{nh}$ (\verb|nh_projection|)
\item tracer advection and biharmonic tracer mixing, or the advection outside the Adams--Bashforth step
\item Adams--Bashforth update of the tracers
\item isoneutral mixing
\item vertical mixing of the tracers
\item density
\item TKE step (\verb|integrate_tke|)
\end{enumerate}

The order of the calls in a time step of the rigid lid and of the implicit free surface is
\begin{enumerate}
\item diffusivities of the TKE closure
\item tracer advection and biharmonic tracer mixing
\item momentum tendencies and hydrostatic pressure (\verb|momentum_tendency_impl|)
\item advective tendency of $w$ (\verb|w_tendency|)
\item forcing of the configuration (\verb|set_forcing|)
\item biharmonic and harmonic friction of $u$, $v$ and of $w$
\item Adams--Bashforth update of $u$, $v$ and of $w$
\item direct forcing of the configuration (\verb|set_forcing_direct|)
\item surface pressure and its gradient (\verb|solve_for_pressure|)
\item vertical friction of $u$, $v$ and of $w$
\item bottom friction
\item surface vertical velocity and the solve for $p_{nh}$ (\verb|nh_projection|)
\item Adams--Bashforth update of the tracers, or the advection outside the Adams--Bashforth step
\item isoneutral mixing
\item vertical mixing of the tracers
\item density
\item TKE step
\end{enumerate}
In both schemes the steps are those of pe\_realistic.
The steps for $w$ are added.
The solve for $p_{nh}$ replaces the vertical velocity from continuity.

The forcing of the configuration may add to the tendency \verb|dw| at the level \verb|tau|, as it does for \verb|du| and \verb|dv|.
The tendency then enters the Adams--Bashforth step.

The solve for $p_{nh}$ comes after the vertical friction and the bottom friction.
The velocity at the end of the time step is therefore free of divergence.
The surface pressure is solved before.
It removes the depth-averaged divergence, which is the mode with the slowest convergence in \eq{eq:pnh}.

The restart needs the fields of pe\_realistic and in addition \verb|dw| and \verb|p_nh|.
Without \verb|p_nh| the solver starts from another first guess.
The run is then no longer identical to the last bit.

\section*{Model parameters and options}

The parameters and options of pe\_realistic are also those of nh\_realistic.
The tables list only the differences.

\subsection*{Required parameters}

\begin{longtable}{@{}llp{8.4cm}@{}}
\hline
Attribute & Unit & Meaning \\
\hline
\endhead
\verb|dxt|, \verb|dyt| & m & widths of the cells in $x$ and $y$, as in pe\_realistic but in metres only \\
\hline
\end{longtable}

\noindent
\verb|coord_degree| must be off.

\subsection*{Options with defaults}

\begin{longtable}{@{}llp{7.6cm}@{}}
\hline
Attribute & Default & Meaning \\
\hline
\endhead
\verb|enable_nonhydrostatic| & on & prognostic $w$ and the solve for $p_{nh}$. Off gives pe\_realistic \\
\verb|epsilon_nh| & $10^{-12}$ & largest divergence (s$^{-1}$) that the solve for $p_{nh}$ leaves \\
\verb|max_iterations_nh| & \verb|max_iterations| & largest number of iterations of the solve for $p_{nh}$ \\
\verb|Ah|, \verb|Ahbi| & unset & as in pe\_realistic, also acting on $w$ \\
\verb|A_v| & $0$ & as in pe\_realistic, also acting on $w$ \\
\verb|weno_advection|, \verb|weno_ref| & off, $1$ & WENO tracer advection as in pe\_realistic.
  It needs \verb|halo_points| $=3$ \\
\hline
\end{longtable}

\noindent
The number of iterations of the last solve for $p_{nh}$ is \verb|congr_itts_nh|.
The divergence it leaves is \verb|estimated_error_nh|.

\subsection*{Snapshot diagnostics}

The snapshot classes \verb|snap_nc4| and \verb|snap_raw| are those of pe\_realistic.
\verb|snap_nc4| also writes $p_{nh}$ under the name \verb|p_nh|.
$w$ is written on the $W$ points as in pe\_realistic.
It is now the prognostic $w$.

\subsection*{Limits}

\begin{itemize}
\item The horizontal coordinates must be Cartesian, see the grid above.
\item The equation for $w$ has no Coriolis term.
\item The energy that the friction of $w$ removes is not given to the TKE closure.
\item With the free surface $w_{surf}$ does not enter the advection of $w$, since the surface face is closed in the equation for $w$. The advection of $w$ then does not conserve $w^2$ exactly in the uppermost $W$ cell.
\item The preconditioner suits $\Delta z \ll \Delta x$. For $\Delta z \approx \Delta x$ the number of iterations grows. A preconditioner with the FFT of a uniform grid, as in Oceananigans, is not available.
\item Runs with different domain decompositions are not identical to the last bit, see the solver above.
\item The LES closures of nh\_simple are not part of nh\_realistic.
\end{itemize}

\end{document}
